\begin{abstract}
Autonomous AI agents now transact at production scale — 69,000 bots executing 165 million transactions across \$50 million USDC in cumulative volume on a single marketplace — without any shared trust layer between participants. Regulatory frameworks (Singapore IMDA, NIST CAISI, EU AI Act) and major AI laboratories (Anthropic, Google) have independently converged on the same structural requirement: an open, portable, cryptographically verifiable trust infrastructure for autonomous agents that no single vendor can deliver alone.

This paper presents MolTrust, a production-deployed implementation of such an infrastructure built on W3C Verifiable Credentials 2.0 and Decentralized Identifiers v1.0, with on-chain anchoring on Base Layer 2. The system architecture is organized around four primitives (identity, authorization, behavioral record, portability), a five-party accountability chain, and the Agent Authorization Envelope (AAE) — a machine-evaluable authorization structure enforced at three layers: cryptographic signatures, API-level credential lifecycle management, and kernel-level syscall monitoring via Falco eBPF integration.

The paper documents three distinguishing capabilities (Claim A): kernel-layer AAE enforcement operating below the agent process boundary; cross-protocol interoperability demonstrated through five reproducible test vectors verified against independent implementations (Claim B); and layered Sybil resistance combining dual-signature interaction proofs, cross-vertical endorsement diversity gating, and principal-DID-linked violation persistence. The reference implementation has been operational since March 2026 across eight credential verticals, with published conformance specifications, nine CWE-mapped security audit checks, and mappings to three independent threat taxonomies.

Empirical validation at adversarial scale is pending. The contribution is deployment-first evidence that the trust infrastructure regulators and industry have converged on is implementable today using W3C-standardized primitives.

\end{abstract}

\section*{Novelty Claims}

This paper presents two contributions:


\textbf{Claim A (Primary):} Production deployment and empirical validation of a W3C Verifiable Credentials and Decentralized Identifier-based trust infrastructure for autonomous agent commerce, distinguished by three specific capabilities: (a) kernel-layer Agent Authorization Envelope enforcement via Falco eBPF integration, operating below the agent process boundary and not bypassable from the agent's own runtime; (b) cross-protocol interoperability through a published conformance specification (CONFORMANCE.md v1.0) with reproducible test vectors (TV-001 through TV-005) verified against independent verifier implementations (qntm Authority Constraints, APS Provider Attestation); and (c) layered Sybil resistance combining dual-signature Interaction Proof Records, cross-vertical endorsement diversity with Jaccard-similarity cluster detection, and principal-DID-linked Violation Records persisting across agent re-registrations.


\textbf{Claim B (Secondary):} Demonstrated protocol-level interoperability through a published conformance specification (CONFORMANCE.md v1.0, SHA-256 drift detection, auto-generated from source), five cross-protocol test vectors (TV-001 through TV-005) verified against qntm Authority Constraints ConstraintEvaluation, and APS Provider Attestation mapping — all reproducible by any party against the live endpoint at api.moltrust.ch.


A companion SSRN preprint presents the formal mapping of the MolTrust Protocol to the AIP feature set identified in Prakash (2026) \cite{ref27_kroehlssrn}; this paper focuses on regulatory alignment, academic positioning, and cross-protocol interoperability.



\section{Introduction}

The deployment of autonomous AI agents has transitioned from research and proof-of-concept into production commerce. Agent.market, launched in April 2026, reports 69,000 autonomous bots executing 165 million transactions with a cumulative volume of \$50 million USDC across seven commercial categories — without any shared trust layer between participants (as reported by the platform at launch) \cite{ref1_agentmarket}. Across financial services, non-human identities now outnumber human employees by 96 to 1; in broader enterprise environments the ratio reaches 144:1 and has grown 44\% in six months \cite{ref2_neville,ref3_entrolabs}. Cloudflare network data covering over 20\% of global web traffic shows AI crawler ratios reaching 30,000 times the equivalent human traffic at Anthropic and 100,000:1 peak crawl-to-refer ratios across leading AI platforms \cite{ref4_prince,ref5_cloudflare}. These figures are not projections. They describe present conditions.


Against this scale, three questions define the agent-to-agent trust gap. They are not new. They are the questions that any counterparty in a human-to-human transaction has answered implicitly for decades through identity documents, credit systems, and institutional reputation. They are what every autonomous agent interaction now confronts from scratch:


\begin{quote}
\textbf{Is this the same agent I interacted with before?}

\textbf{Is this agent permitted to do what it claims to do?}

\textbf{Has this agent acted consistently with its declared parameters in past interactions?}

\end{quote}

These three questions map to distinct primitives — identity, authorization, behavioral history — and they are what the Singapore IMDA Framework, the NIST NCCoE Concept Paper, and Anthropic's \textit{Trustworthy Agents in Practice} paper independently identify as the present gap in production infrastructure. Identity infrastructure designed for human users (OAuth 2.0, SAML, LDAP) and for static service accounts does not survive agent cloning, redeployment, or cross-platform migration. AI framework-native trust mechanisms address the problem within single platform boundaries but not across them. W3C Decentralized Identifiers and Verifiable Credentials provide correct primitives \cite{ref6_w3cdid,ref7_w3cvc} but require application-layer specificity for agent commerce workflows.


Three distinct signals from mid-2025 through early 2026 indicate that this gap is now explicitly recognized by regulators, institutional standards bodies, and major AI laboratories.


On 22 January 2026 at the World Economic Forum, the Singapore Infocomm Media Development Authority (IMDA) released the \textit{Model AI Governance Framework for Agentic AI} — the first comprehensive regulatory framework for autonomous AI systems \cite{ref8_imda}. The Framework's technical controls dimension explicitly requires "strong authentication measures such as scoped API keys, per-agent identity tokens, and robust observability such as the logging of tool calls and access history" [8, p. 14]. On 17 February 2026, the U.S. National Institute of Standards and Technology (NIST) launched the AI Agent Standards Initiative through its Center for AI Standards and Innovation (CAISI), framing agent identity and authorization as a national-priority standards gap \cite{ref9_nistcaisi}. The associated NCCoE concept paper, \textit{Accelerating the Adoption of Software and AI Agent Identity and Authorization} (February 2026), directly addresses the enterprise security gap where AI agents are commonly treated as generic service accounts without dedicated identity, authorization, or accountability controls \cite{ref10_nccoe}. On 9 April 2026, Anthropic published \textit{Trustworthy Agents in Practice}, stating explicitly that the model layer alone cannot secure agentic AI and calling on the industry to build "shared infrastructure that no single company can deliver alone" \cite{ref11_anthropic}. Google's Secure AI Framework, extended to SAIF 2.0 with a dedicated Agent Risk Map in 2026, identifies the same architectural gap and donated its risk map data to the Coalition for Secure AI \cite{ref12_google}.


These signals converge on a structural conclusion: an open, portable, cryptographically verifiable trust layer for autonomous agents is not a research proposal. It is a regulatory and operational requirement that the industry has collectively acknowledged but has not yet collectively delivered.


This paper describes MolTrust, a production-deployed implementation of such an infrastructure. The MolTrust Protocol is built on W3C Verifiable Credentials 2.0 and W3C Decentralized Identifier Core 1.0, with all credential issuance and violation records anchored on Base Layer 2. The reference implementation is live at api.moltrust.ch, has been operational since March 2026, and issues credentials across eight verticals (Core Identity, Commerce, Travel, Skill Verification, Prediction Markets, Brand Protection, Music, Sports Integrity). An embedded ERC-8004 ecosystem scanner indexes 44,355 agents observed across the on-chain agent registry, providing cross-ecosystem visibility independent of MolTrust registration status. The Agent Authorization Envelope (AAE) — a machine-evaluable structure encoding MANDATE, CONSTRAINTS, and VALIDITY blocks — enforces authorization boundaries at three architectural layers: cryptographic signatures (Ed25519 + RFC 8785 JCS), API-level trust score and credential lifecycle management, and kernel-level syscall monitoring via Falco eBPF integration \cite{ref13_falcobridge}. As of April 2026, all three enforcement layers are operational.


The contribution of this work is not the invention of new cryptographic primitives or identity concepts. The contribution is \textbf{the demonstration that the primitives standardized by W3C, combined with on-chain anchoring and kernel-layer enforcement, are sufficient to build the trust infrastructure that NIST, IMDA, and major AI laboratories identify as missing — and that such infrastructure operates today in production, across eight credential verticals, with kernel-layer AAE enforcement, reproducible cross-protocol conformance evidence against independent verifier implementations, and layered Sybil resistance. Each of these distinguishing mechanisms is independently verifiable against the live endpoint and the published conformance specification.} This deployment-first contribution is complementary to, and validates, a growing body of academic literature on multi-agent security \cite{ref14_schroederdewitt}, inter-agent trust models \cite{ref15_hurong}, governance architectures \cite{ref16_syros,ref17_acharya}, conceptual DID-and-VC frameworks for agents \cite{ref36_rodriguezgarzon}, and threat taxonomies \cite{ref18_ferrag,ref19_owasp,ref20_mao}.


The remainder of this paper is organized as follows. Section 2 positions MolTrust within the regulatory context (NIST, IMDA, EU AI Act), the industry framework convergence (Anthropic, Google, Microsoft), and the academic literature. Section 3 describes the MolTrust system architecture: the four primitives, the Five-Party Trust Chain, the Agent Authorization Envelope, and the three-layer architecture. Section 4 presents deployment and empirical evidence, including the published conformance specification (CONFORMANCE.md v1.0) with explicit mapping to multiple independent threat taxonomies. Section 5 describes cross-protocol interoperability — the basis of the secondary novelty claim. Section 6 addresses Sybil resistance and behavioral consistency. Section 7 discusses limitations and compares MolTrust to the closest adjacent academic frameworks. Section 8 concludes.



\section{Related Work and Positioning}

The MolTrust Protocol occupies a specific position across four bodies of work: regulatory and institutional standards that define requirements for agent trust infrastructure (Section 2.1); industry frameworks from major AI laboratories and platform vendors that describe the problem and propose partial solutions (Section 2.2); academic literature on agent trust architectures and trust-model taxonomies (Section 2.3); and academic literature on threat taxonomies and closely adjacent theoretical frameworks (Section 2.4).


\subsection{The Regulatory and Institutional Context}

\textbf{Singapore IMDA Model AI Governance Framework for Agentic AI.} Published 22 January 2026 at the World Economic Forum, the IMDA Framework is the world's first comprehensive governance framework specifically for agentic AI — AI systems capable of autonomous planning, reasoning, and action \cite{ref8_imda}. The Framework is organized around four dimensions: (1) \textit{Assess and Bound the Risks Upfront} through use-case selection, agent limits, and least-privilege access; (2) \textit{Make Humans Meaningfully Accountable} through defined checkpoints requiring human approval; (3) \textit{Implement Technical Controls and Processes} across the agent lifecycle; and (4) \textit{Enable End-User Responsibility} through transparency and training. The Framework states that each AI agent should have a verifiable digital identity, enabling a clear audit trail of which agent acted, under whose authorisation, and when — and requires strong authentication measures including scoped API keys, per-agent identity tokens, and robust observability such as the logging of tool calls and access history \cite{ref8_imda,ref21_imdafull}. The MolTrust Protocol implements each of these four dimensions directly (detailed mapping in Section 4).


\textbf{NIST AI Agent Standards Initiative and NCCoE Concept Paper.} On 17 February 2026, NIST's Center for AI Standards and Innovation (CAISI) launched the AI Agent Standards Initiative \cite{ref9_nistcaisi}, organized around three pillars: industry-led standards development with U.S. leadership in international bodies, community-led open-source protocol development, and fundamental research in agent security and identity infrastructure. The associated National Cybersecurity Center of Excellence (NCCoE) concept paper, \textit{Accelerating the Adoption of Software and AI Agent Identity and Authorization} (5 February 2026), explicitly frames the gap that MolTrust addresses: AI agents are commonly treated as generic service accounts without dedicated identity, authorization, or accountability controls \cite{ref10_nccoe}. The paper identifies OAuth 2.0, OpenID Connect, and SPIFFE/SPIRE as foundation identity standards that must be extended for autonomous agents. The NIST Control Overlays for Securing AI Systems (COSAiS) project is developing SP 800-53 control overlays for single-agent and multi-agent deployment scenarios \cite{ref22_cosais}. As of this writing, the COSAiS overlays have not been finalized; MolTrust's implementation of NIST AI RMF controls is documented in a separate methodology note.


\textbf{EU AI Act.} Regulation (EU) 2024/1689 establishes transparency and risk classification requirements for AI systems with enforcement beginning August 2026 \cite{ref23_euaiact}. Article 14 (Human Oversight) and Article 15 (Accuracy, Robustness and Cybersecurity) apply directly to autonomous agents deployed in high-risk domains. MolTrust's AAE structure enforces Article 14 requirements through the delegated\_authority and approval threshold mechanisms; Article 15 requirements are addressed through the Skill Audit infrastructure (nine conformance checks including a secrets scan mapped to CWE-798) and through the Violation Record mechanism with on-chain anchoring.


\textbf{Convergence signal.} The IMDA Framework, NIST CAISI Initiative, and EU AI Act — issued independently across three jurisdictions within an eleven-month window — identify substantively the same requirements: verifiable agent identity, machine-readable authorization boundaries, behavioral accountability, and human oversight checkpoints. These requirements are what MolTrust implements. They are also what the industry frameworks in Section 2.2 identify as the missing layer.


\subsection{Industry Framework Convergence}

\textbf{Anthropic: Trustworthy Agents in Practice.} On 9 April 2026, Anthropic published its framework for trustworthy agent development, organized around five principles: human control, alignment with human values, securing interactions, transparency, and privacy \cite{ref11_anthropic}. The framework explicitly states that the model layer alone cannot secure agentic AI and that the industry requires shared infrastructure that no single vendor can provide unilaterally. Anthropic identifies four components determining agent behavior — model, harness (instructions and guardrails), tools, and environment — and argues that a well-trained model can still be exploited through overly permissive tools or poorly configured environments. The framework explicitly calls for open standards over proprietary alternatives and cites Anthropic's donation of the Model Context Protocol to the Linux Foundation as an example of this approach. MolTrust's position is direct: the W3C VC + DID infrastructure is an open, standards-based realization of the shared trust layer Anthropic's framework describes as necessary. MolTrust's AAE provides the structural separation between agent capabilities (what the model can be asked to do) and agent authorization (what the agent is permitted to do), enforced cryptographically rather than through model alignment alone.


\textbf{Google DeepMind: Secure AI Framework (SAIF) 2.0 and Agent Risk Map.} Google's SAIF 2.0, extended in 2026 with a dedicated Agent Risk Map, identifies fifteen common AI risks and maps them to specific controls across four component areas: Data, Infrastructure, Model, and Application \cite{ref12_google}. Agent-specific risks identified include Rogue Actions (expanded action capability without corresponding alignment controls), Insecure Integrated Components (vulnerabilities in tool and plugin integrations), Prompt Injection (confusion between instructions and data), and Sensitive Data Disclosure (agent-privileged access to user information). Google's three core principles for agent security — well-defined human controllers, carefully limited powers, and observable actions and planning — align with IMDA Framework Dimensions 2, 1, and 3 respectively. Google donated SAIF risk-map data to the Coalition for Secure AI (CoSAI), whose founding members include Anthropic, Cisco, IBM, Intel, Nvidia, and PayPal \cite{ref24_cosai}. MolTrust's AAE MANDATE and CONSTRAINTS blocks provide a machine-evaluable realization of Google's "carefully limited powers" principle; the Interaction Proof Record infrastructure provides a cryptographically verifiable realization of "observable actions and planning."


\textbf{Microsoft: Entra Agent ID.} Microsoft announced Entra Agent ID at Ignite 2025 as an identity and access management framework extending Microsoft Entra capabilities to AI agents \cite{ref25_entra}. The platform provides identity constructs for AI agents integrated with OAuth 2.0, MCP protocol support, Conditional Access policies, and lifecycle governance within the Microsoft Entra ecosystem. Entra Agent ID addresses agent identity comprehensively for enterprises operating within the Microsoft identity perimeter. The MolTrust Protocol addresses a different scope: cross-organizational agent-to-agent trust based on W3C Verifiable Credentials and Decentralized Identifiers, verifiable by any party independently of the issuer's platform. The two approaches operate at different layers of the agent identity stack and are not substitutive.


\textbf{Convergence signal.} The Anthropic, Google, and Microsoft frameworks — issued between October 2025 and April 2026 — independently identify agent identity, authorization granularity, and behavioral observability as the three axes of agent security. They differ in implementation approach: Anthropic calls for open shared infrastructure, Google contributes its framework to industry consortia, and Microsoft provides a platform-native solution. MolTrust contributes an open-standards, vendor-independent realization of the shared layer Anthropic describes and Google's principles require, designed from the outset for cross-organizational verification.


\subsection{Academic Foundations — Architecture and Taxonomy}

\textbf{Syros et al., SAGA (NDSS 2026).} Syros, Suri, Ginesin, Nita-Rotaru, and Oprea of Northeastern University Khoury College proposed SAGA (Security Architecture for Governing AI Agentic Systems), accepted for NDSS 2026 \cite{ref16_syros}. SAGA introduces a centralized Provider architecture: users register with a Provider entity that maintains agent contact information and user-defined access control policies, issues cryptographic One-Time Keys (OTKs) for inter-agent communication, and enforces policies at the token level. The architecture provides formal security guarantees and empirical evaluation across multiple agentic tasks with minimal performance overhead.


MolTrust and SAGA address the same problem — secure and accountable agent-to-agent interaction — with structurally different architectures. SAGA is centralized: a single Provider is required, and all agent identity resolution passes through that Provider. MolTrust is decentralized: agent DIDs are cryptographically self-verifying through W3C-standardized resolution, and any party holding an issuer's DID Document can verify credentials without contacting a central Provider. The trade-off is explicit. SAGA achieves tighter policy enforcement through centralization; MolTrust achieves portability and jurisdictional independence through decentralization. The two architectures are not interchangeable and they are not competing for identical use cases. SAGA is optimal for organizations deploying agents under a single administrative domain where a Provider can plausibly be trusted by all participants. MolTrust is optimal for agent-to-agent commerce across organizational and jurisdictional boundaries where no such Provider exists.


\textbf{Hu and Rong, Inter-Agent Trust Models (AAAI 2026 TrustAgent Workshop).} Botao "Amber" Hu (DPhil candidate, University of Oxford Department of Computer Science) and Helena Rong (Assistant Professor, NYU Shanghai) proposed a six-category taxonomy of inter-agent trust models: Brief (self- or third-party verifiable claims), Claim (self-proclaimed capabilities without external verification), Proof (cryptographic verification including ZKP and TEE attestations), Stake (bonded collateral with slashing), Reputation (crowd feedback and graph-based trust signals), and Constraint (sandboxing and capability bounding) \cite{ref15_hurong}. The paper analyzes how A2A, AP2, and ERC-8004 protocols map to these categories and argues for hybrid approaches combining multiple trust signals to resist manipulation.


MolTrust implements a quad-layer combination of Hu and Rong's categories: \textbf{Brief} (W3C VCs issued by third-party authorities including Trust Tier 0 KYC-backed developer credentials), \textbf{Proof} (Ed25519 signatures over canonical JSON per RFC 8785, with on-chain anchoring on Base L2), \textbf{Reputation} (Trust Score derived from endorsement graph with cross-vertical diversity and Jaccard similarity sybil detection), and \textbf{Constraint} (AAE MANDATE and CONSTRAINTS blocks with default-deny semantics and attenuation-only delegation). Stake is supported as optional for agent registration but is not a primary signal. This multi-layer architecture is consistent with Hu and Rong's argument that no single trust mechanism is robust against sophisticated adversaries and that hybrid signal combination is a structural requirement rather than an implementation optimization.


\textbf{Schroeder de Witt, Open Challenges in Multi-Agent Security (Oxford, 2025).} Christian Schroeder de Witt (University of Oxford Department of Engineering Science) introduced multi-agent security as a distinct research field at the NeurIPS 2023 workshop and formalized the field's scope in a 2025 position paper \cite{ref14_schroederdewitt}. The paper identifies threats that emerge specifically from agent interactions: secret collusion through steganographic communication \cite{ref26_motwani}, coordinated swarm attacks, network-propagated privacy breaches and jailbreaks, and dispersion/stealth optimization by adversarial agents. These threats are distinct from traditional cybersecurity threats (which target individual systems) and from AI safety threats (which target individual model behavior).


The MolTrust Protocol addresses the multi-agent security threat landscape in three structural ways. First, the Interaction Proof Record infrastructure produces cryptographically verifiable records of specific interactions between identified agents, making collusion empirically auditable rather than detectable only through behavioral inference. Second, the dual-signature requirement for bilateral interaction proofs (Section 6) prevents single-adversary fabrication of interaction histories, constraining the scale at which a single actor can manipulate the behavioral record. Third, the cross-vertical diversity requirement in the Reference Reputation Model (Section 3.1) and the Jaccard similarity sybil detection heuristic address the coordinated-swarm and dispersion threats Schroeder de Witt identifies, though explicitly as lightweight heuristics rather than robust adversarial defenses.


\subsection{Academic Foundations — Threats and Closest Parallels}

\textbf{Acharya, Trustless Intent Verification for Autonomous Agents (TIVA, 2025).} Vivek Acharya proposed TIVA, a blockchain-based framework for verifying the authenticity and intent of AI-initiated payments \cite{ref17_acharya}. TIVA combines decentralized identity (DIDs and verifiable credentials) for agent authentication, on-chain intent proofs capturing user authorization, zero-knowledge proofs for privacy-preserving policy compliance, and optional TEE attestations for agent execution integrity.


\textbf{Rodriguez Garzon et al., AI Agents with Decentralized Identifiers and Verifiable Credentials (TU Berlin, 2025).} Rodriguez Garzon and colleagues present a closely related conceptual framework in which each AI agent is endowed with a self-sovereign digital identity combining a ledger-anchored W3C DID with a set of third-party issued W3C VCs \cite{ref36_rodriguezgarzon}. Their prototypical multi-agent implementation demonstrates technical feasibility of cross-domain trust establishment between agents through the spontaneous exchange of DID-bound VCs. The paper identifies important limitations when the agent's LLM alone is in sole charge of security procedures — a finding that directly motivates enforcement architectures that operate independently of the model layer.


TIVA and Rodriguez Garzon et al. are the closest conceptual parallels to MolTrust in the academic literature. The architectural primitives are substantively identical across all three: DIDs + VCs for identity, on-chain anchoring for accountability, cryptographic signatures for authorization. The frameworks are complementary in focus: Acharya targets payment-intent verification; Rodriguez Garzon et al. target the broader trust-establishment dialogue at the onset of agent-to-agent interaction. MolTrust operationalizes both, with the AAE VALIDITY block capturing the trust-establishment binding Rodriguez Garzon et al. describe and the AAE CONSTRAINTS.limits block capturing the payment-intent boundaries Acharya formalizes. The difference between these academic frameworks and the present work is in deployment status and enforcement scope: both cited papers describe conceptual frameworks with prototypical implementations \cite{ref17_acharya,ref36_rodriguezgarzon}. MolTrust is a production-deployed system with a live API, published conformance specifications with reproducible test vectors, and kernel-layer enforcement operating below the agent process boundary. The contribution of this paper is therefore complementary: MolTrust demonstrates that the architectural pattern these academic frameworks describe is deployable today, at operational scale, with independently verifiable evidence — and with enforcement mechanisms (Layer 3 kernel monitoring) that address precisely the "LLM-in-sole-charge" limitation Rodriguez Garzon et al. identify.


\textbf{Ferrag et al., From Prompt Injections to Protocol Exploits (ICT Express 2026).} Ferrag, Tihanyi, Hamouda, Maglaras, Lakas, and Debbah, affiliated with Guelma University (Algeria), Technology Innovation Institute (Abu Dhabi), and Edinburgh Napier University, published a comprehensive threat taxonomy of LLM-agent ecosystems in ICT Express (Elsevier, in press 2026, DOI: 10.1016/j.icte.2025.12.001) \cite{ref18_ferrag}. The taxonomy catalogues more than thirty attack techniques across four domains: Input Manipulation (prompt injections, long-context hijacks, multimodal adversarial inputs), Model Compromise (backdoors, poisoning), System and Privacy Attacks (side-channels, membership inference, retrieval poisoning), and Protocol Vulnerabilities (exploits in MCP, ACP, ANP, and A2A). The framework is validated through expert review and cross-mapping with CVE and NIST National Vulnerability Database entries — a direct bridge to U.S. institutional vulnerability tracking.


MolTrust's coverage mapping against Ferrag et al.'s taxonomy is detailed in Section 4. The Protocol Vulnerabilities domain is directly in MolTrust's scope; Input Manipulation and Model Compromise domains are outside the scope of a trust infrastructure layer and remain the responsibility of the deploying organization. The CVE/NIST NVD cross-mapping methodology Ferrag et al. establish is the reference approach for the MolTrust NIST AI RMF mapping methodology note.


\textbf{Mao et al., Systematization of Knowledge on AI Agent Security (2026).} Mao et al. published a systematization of knowledge identifying five dimensions of AI agent security — Agent Integrity, Transaction Authorization, Inter-Agent Trust, Market Manipulation, and Regulatory Compliance — and twelve cross-layer attack vectors \cite{ref20_mao}. The systematization provides a useful cross-check: it is produced independently from the western academic work discussed above and arrives at substantively convergent conclusions about the threat landscape. MolTrust addresses 8 of Mao et al.'s 12 cross-layer attack vectors through live capabilities; 2 vectors are on the published roadmap (N2C cumulative exposure monitoring Phase 2, R2I AML pattern detection Phase 2/3); 2 vectors are explicit non-coverage (M2A model integrity via Sigstore partnership space; O2P oracle integrity via Chainlink partnership space). The full coverage matrix is presented in Section 4. Mao et al.'s work is referenced in this paper as an independent confirming systematization, not as the primary threat framework; the primary threat references are Ferrag et al. (peer-reviewed journal publication with CVE/NIST NVD integration), Schroeder de Witt (Oxford field-defining position paper), and the convergent regulatory documents (NIST, IMDA, EU AI Act) discussed in Section 2.1.


\subsection{Position Summary}

MolTrust occupies a specific position across these four bodies of work. It implements the technical controls that NIST CAISI, Singapore IMDA, and the EU AI Act require. It provides an open, vendor-independent realization of the shared infrastructure Anthropic identifies as necessary and that Google's SAIF 2.0 principles describe. It offers a cross-organizational, W3C-standardized trust layer that operates at a different scope than the platform-native agent identity solutions deployed by major cloud vendors. It is the deployed realization of the architectural pattern Acharya describes theoretically as TIVA. It operationalizes the quad-layer trust-signal combination (Brief + Proof + Reputation + Constraint) that Hu and Rong argue is structurally required. And it addresses the multi-agent security threats Schroeder de Witt identifies, through mechanisms mapped to the Ferrag et al. taxonomy with convergent confirmation from the Mao et al. systematization.


The contribution of this paper is not a new framework. It is empirical evidence, from a live deployment, that the framework the regulators, institutions, and industry have converged on is implementable today, across multiple credential verticals, with reproducible conformance test vectors, using W3C-standardized primitives and on-chain anchoring. Section 3 describes the system architecture.



\section{System Architecture}

The MolTrust Protocol is organized around four primitives, a five-party trust chain, the Agent Authorization Envelope, and a three-layer enforcement architecture. This section defines each in a form suitable for academic review. Complete normative definitions, data schemas, and signing rules are published in the MolTrust Protocol Technical Specification v0.8 \cite{ref28_techspec}; this paper presents the architecture at the level of detail appropriate for positioning within the literature.


\subsection{Four Primitives}

The MolTrust Protocol is constructed from four primitives, each addressing one of the trust questions identified in Section 1.


\textbf{Identity.} Every agent in the protocol is assigned a Decentralized Identifier (DID) conforming to W3C DID Core v1.0 \cite{ref6_w3cdid}. DIDs are globally unique, cryptographically verifiable, and require no central registry to create or verify. Ownership is proven by possession of the corresponding Ed25519 private key. The reference DID method is \texttt{did:moltrust}, registered and operated by the MolTrust reference registry; any conformant implementation MAY use any DID method supporting Ed25519 or secp256k1 signing. A DID provides a stable, portable reference to an agent across platforms and over time — it survives redeployment, migration, and platform changes. This answers the first trust question: \textit{is this the same agent I interacted with before?}


\textbf{Authorization.} The relationship between an agent and the principal on whose behalf it acts is expressed as a W3C Verifiable Credential (VC) conforming to the Verifiable Credentials Data Model 2.0 \cite{ref7_w3cvc}. Credentials are digitally signed, tamper-evident attestations issued by the granting party, describing what an agent is permitted to do, under what conditions, and until when. Credential chains support delegation: an organization may authorize an agent, which in turn authorizes a sub-agent, with each step cryptographically linked to the previous. Any verifier can traverse the chain independently without contacting the original issuer. This answers the second trust question: \textit{is this agent permitted to act as it claims?}


\textbf{Behavioral Record.} Interactions between agents produce Interaction Proof Records (IPRs) — cryptographic records that a specific interaction occurred at a specific time between identified parties, with a recorded outcome. Bilateral IPRs carry sequential dual signatures: the initiator signs first, the responder's signature covers the initiator's signature, creating a non-repudiable commitment chain. Completed IPRs are Merkle batch-anchored on Base Layer 2. IPRs accumulate into a Trust Score reflecting observed behavior across time and context.


The Trust Score is computed by the registry from four weighted components plus additive modifiers:


\begin{lstlisting}[basicstyle=\ttfamily\small]
raw = 0.6 x direct_score
    + 0.3 x propagated_score
    + 0.1 x cross_vertical_bonus
    + interaction_bonus              # additive, capped

final = clamp(raw - (sybil_penalty x 20) + inactivity_penalty, 0, 100)
         (seeds: final = max(final, base_score))
\end{lstlisting}

The \textbf{direct-endorsement score} (weight $\alpha$ = 0.6) is derived from endorsements the agent has received, weighted by each endorser's own trust score and time-decayed exponentially (half-life 90 days). The \textbf{propagated endorsement score} (weight $\beta$ = 0.3) reflects second-order endorser reputation. The \textbf{interaction bonus} is additive (un-coefficiented) and capped at 10 points, derived from verified Interaction Proof Records via \texttt{compute\_ipr\_bonus}. The \textbf{inactivity penalty} is an additive modifier applied before the final clamp operation, penalizing agents with prolonged absence of verifiable activity (addresses RSA Conference 2026 Agent Security Gap Analysis, Gap 3; full mechanics in Technical Specification v0.8 \cite{ref28_techspec}).


\textbf{Cross-vertical diversity} is implemented as a two-stage mechanism. The \textbf{incentive layer} applies a weighted bonus ($\gamma$ = 0.1) capped at 30 points (min(unique\_verticals $\times$ 10, 30)), rewarding breadth of trust attestation across distinct verticals. The \textbf{gate layer} applies a flat Sybil penalty of 10.0 when a non-seed agent has endorsements from fewer than 3 distinct verticals, producing effective score nullification after the $\times$20 multiplier (10.0 $\times$ 20 = 200, clamped to [0, 100]).


The three-vertical threshold reflects a deliberate design choice against Sybil coordination. Two colluding agents can fabricate mutual endorsements within a single vertical at minimal cost; requiring endorsements from three independent verticals increases coordination cost non-linearly, as each additional vertical requires the Sybil operator to maintain distinct credibility signals across unrelated domains. The threshold is calibrated against the seed bootstrap mechanism described below and is empirically monitored; adjustment pathways are preserved via configuration.


"Distinct verticals" refers to the canonical vertical categories tracked by the registry: \texttt{core}, \texttt{skill}, \texttt{shopping}, \texttt{travel}, \texttt{prediction}, \texttt{salesguard}, \texttt{sports}, plus eight credential-type verticals (\texttt{VerifiedSkillCredential}, \texttt{BuyerAgentCredential}, \texttt{AuthorizedAgentCredential}, \texttt{TravelAgentCredential}, \texttt{PredictionTrackCredential}, \texttt{ProductProvenanceCredential}, \texttt{AuthorizedResellerCredential}, \texttt{SkillEndorsementCredential}). Unique vertical count includes both endorsement-declared verticals and credential-type verticals.


The \textbf{Sybil penalty} is proportional to the Jaccard similarity of the agent's endorser set with its most similar peer, triggered when similarity exceeds 0.8 (production threshold). The final score is clamped to [0, 100]. Scores are withheld until a non-seed agent has endorsements from at least three distinct endorser DIDs.


The gate mechanism necessitates a bootstrap path. \textbf{Seed agents}, established during initial registry deployment, are protected by a seed-floor-guard: their trust score cannot fall below their assigned \texttt{base\_score} regardless of penalty accumulation. Seed base\_scores are calibrated against the gate's severity and intentionally below the theoretical ceiling: TrustScout (85.0), MolTrust Ambassador (80.0), VCOne (75.0), seeded agent (70.0), AgentNexus (60.0). This tiering reflects a transparent bootstrap hierarchy rather than privileged status; all seed floors and their base\_scores are discoverable via the registry's \texttt{/swarm/stats} endpoint.


The reference formula is specified normatively in Section 4 of the Technical Specification \cite{ref28_techspec}; this paper reproduces the structural components to allow independent readers to evaluate the model without consulting the normative specification.


This answers the third trust question: \textit{has this agent acted consistently with its declared parameters?}


\textbf{Portability.} All three primitives above are expressed in open, platform-independent formats. A credential issued by one organization is verifiable by any counterparty. A behavioral record compiled from interactions on one platform travels with the agent to any other. Portability in this sense means interoperability of evidence formats, not identity of scoring outcomes. Different implementations consuming the same evidence may produce different trust scores depending on their scoring model. The protocol standardizes \textit{what} evidence is recorded and \textit{how} it is structured — not how it is ultimately weighted. This property is what distinguishes MolTrust from platform-scoped agent identity solutions, and what makes it suitable for the cross-organizational agent commerce scenario Section 2.2 describes as the structural gap.


\subsection{The Five-Party Trust Chain}

Trust in the autonomous agent economy does not begin with an agent. It begins with the humans and organizations that deploy them. The MolTrust Protocol formalizes a five-party accountability chain:


\textbf{Developer $\rightarrow$ Owner $\rightarrow$ Agent $\rightarrow$ Instructor $\rightarrow$ Counterparty}


Each link is cryptographically verifiable. The Developer is the entity that wrote or deployed the agent code; the Developer MAY hold a Trust Tier 0 KYC-backed DID binding the agent's authorship to a verified real-world identity. The Owner is the organization or individual operating the agent (potentially different from the Developer); ownership transfers are recorded immutably. The Agent is the autonomous system acting on behalf of the Owner; the Agent holds its own DID and VCs and cannot shed identity through redeployment. The Instructor is the entity providing runtime instructions (human, agent, or workflow); Interaction Proofs bind instructions to outcomes. The Counterparty is the system or agent the Agent interacts with; the Counterparty can verify the Agent's credentials independently without calling the MolTrust API.


A key structural property: Agent DIDs are independent from Developer DIDs. If a Developer credential is revoked — for any reason, including KYC-provider compromise, jurisdictional dispute, or accidental exposure — the agents operated under that Developer identity are \textit{not} automatically invalidated. Each agent's credential chain is evaluated independently, and each agent's Trust Score reflects its own behavioral record. This decoupling is deliberate: it prevents a single point of failure in the human accountability layer from cascading into operational disruption of the agent layer, while preserving the traceability that regulatory frameworks require.


\subsection{Agent Authorization Envelope (AAE)}

Every MolTrust credential may carry, or be accompanied by, a machine-evaluable authorization object — the Agent Authorization Envelope (AAE). The AAE is organized in three normative blocks: \textbf{MANDATE} (what is the agent permitted to do?), \textbf{CONSTRAINTS} (under what conditions?), and \textbf{VALIDITY} (is this authorization still effective?).


\textbf{MANDATE} defines the agent's permitted purpose (enumerated categories: \texttt{commerce}, \texttt{data\_read}, \texttt{data\_write}, \texttt{communication}, \texttt{delegation}, \texttt{administration}), the URI patterns of allowed actions, explicitly denied actions (which take precedence over allowed actions), resource patterns for attribute-based access control, and delegation sub-fields controlling whether and how the agent may authorize sub-agents.


\textbf{CONSTRAINTS} defines operational boundaries: time windows (time-to-live, session duration, allowed days and hours with timezone), financial thresholds (autonomous threshold below which no additional check is required, step-up threshold triggering additional verification, approval threshold requiring human sign-off, maximum transactions per hour, specified currency), jurisdictions (ISO 3166-1 alpha-2 country codes), counterparty minimum trust score, and obligations (tool allowlists, human approval thresholds for specific action types).


\textbf{VALIDITY} defines issuer identity, holder binding (the DID to which this envelope is bound — the envelope MUST NOT be used by any other agent), issuance and expiry timestamps (open-ended credentials are not permitted), revocation endpoint URL, and optional on-chain anchor reference.


The AAE implements four validation rules that any conformant evaluator MUST enforce: \textbf{default-deny} (any action not explicitly matched by an allowed pattern is denied), \textbf{deny-precedence} (explicit denial overrides any allow), \textbf{attenuation-only delegation} (a delegated AAE must be a strict subset of its parent), and \textbf{mandatory expiry} (no AAE may be open-ended).


The design of the AAE draws on established specifications. The constraint model is informed by NIST SP 800-162 (Guide to Attribute-Based Access Control) \cite{ref29_abac}. The authorization request semantics are consistent with RFC 9396 (Rich Authorization Requests) \cite{ref30_rar}. The canonical serialization for signing is RFC 8785 (JSON Canonicalization Scheme) \cite{ref31_jcs}. The signature algorithm is Ed25519 per Ed25519Signature2020 \cite{ref32_ed25519}. These standards choices are deliberate: they are the most widely implemented, most reviewed, and most permissively licensed options available, which is consistent with the MolTrust design principle that the reference implementation should be replaceable by any conformant alternative.


\subsection{Three-Layer Enforcement Architecture}

The Agent Authorization Envelope defines what an agent may do. Definition alone is insufficient — constraints must be enforced. MolTrust implements a three-layer enforcement model.


\textbf{Layer 1 — Cryptographic.} Ed25519 signatures and RFC 8785 canonical JSON canonicalization make AAE boundaries tamper-proof by construction. Any counterparty can verify that an AAE has not been modified since issuance, without infrastructure dependency. This is the baseline enforcement layer and is active in every verification.


\textbf{Layer 2 — API.} The MolTrust registry enforces constraints at the application level: Trust Score degradation on violation, Interaction Proof Record submission for audit trails, credential revocation for compromised agents, CAEP-compatible revocation event propagation. Revocation propagates to verifiers within 60 seconds in the reference implementation via cache invalidation. Verifier-side behavior on unreachable revocation endpoints defaults to fail-closed (credential denied) in line with the operational threat model for autonomous agent commerce.


\textbf{Layer 3 — Kernel.} Falco eBPF integration monitors agent behavior at the Linux syscall level \cite{ref13_falcobridge}. If an AAE declares that an agent may not write to a specific resource pattern, Falco detects the syscall attempt and submits a violation proof to MolTrust — regardless of whether the agent's own runtime reports it. This is enforcement \textit{below} the agent's process boundary. The agent cannot suppress or modify the detection signal from userspace.


Layer 3 is a qualitative difference from Layers 1 and 2. Cryptographic and API enforcement rely on the agent's environment to report violations; kernel enforcement operates independently of the agent. The infrastructure itself enforces the constraint, and any policy violation attempt — even one ultimately blocked by container permissions — becomes a permanent, verifiable event in the agent's trust history. As of April 2026, the Falco Bridge is live on the OpenClaw gateway instance. Reference implementation source is published at \texttt{github.\allowbreak com/\allowbreak HaraldeRoessler/\allowbreak moltrust-falco-bridge}.


The three layers are complementary, not redundant: \textbf{cryptography} proves identity and authorization at issuance time; \textbf{API verification} proves compliance at request time; \textbf{kernel enforcement} proves behavioral integrity at runtime. No single layer is sufficient. Together they close the gap that industry analysis identifies as the hardest open problem in agent identity: detection of policy self-modification between authorization and execution.


\subsection{Governance Boundaries and Decentralization Roadmap}

An honest account of the MolTrust architecture must distinguish between decentralization at the \textit{evidence layer} and centralization at the \textit{registry layer}. The evidence produced by MolTrust — DID Documents, Verifiable Credentials, Interaction Proof Records, Violation Records — is cryptographically self-verifying: any party holding an issuer's public key can validate any signed artifact without consulting any central service. In this sense, the protocol is genuinely decentralized and verifier-independent. The \texttt{did:moltrust} method and the reference registry at api.moltrust.ch, however, are operated by a single organization (CryptoKRI GmbH). This is intentional for the reference implementation; it is not a claim of structural decentralization at the registry layer.


Three forms of decentralization are relevant, and MolTrust's current position on each is explicit. \textbf{Method-level decentralization:} any DID method supporting Ed25519 or secp256k1 signing can produce MolTrust-compatible credentials — the reference method is one of several conformant options, not the only path. \textbf{Registry-level decentralization:} independent operators may run conformant registries using their own infrastructure, with their own DID method namespaces; the Technical Specification \cite{ref28_techspec} defines Layer B conformance requirements for such independent registries. \textbf{Governance-level decentralization:} protocol evolution is currently governed by a single operator without a formal multi-stakeholder process. A W3C DID Method Registry submission is planned for Protocol v0.9, which will subject the \texttt{did:moltrust} method to external specification review. A broader community governance model, including a conformance certification process independent of the reference implementation operator, is identified as necessary roadmap work but is not yet in place. The protocol is therefore partially decentralized by design and partially centralized by current operational reality; the roadmap is for the operational reality to converge toward the design.


This distinction matters for academic review: claims of decentralization at the evidence layer are verifiable and accurate; claims of decentralization at the governance layer are aspirational and are explicitly labeled as such. The NIST NCCoE concept paper \cite{ref10_nccoe} identifies similar tension in any identity system that combines open standards with operator-specific deployment; the MolTrust position is to acknowledge the tension rather than resolve it rhetorically.



\section{Deployment and Empirical Evidence}

This section presents the evidence underlying Novelty Claim A: that MolTrust is a production-deployed implementation of the trust infrastructure Section 2 identifies as the convergent regulatory, industry, and academic requirement. Subsection 4.1 summarizes the live deployment. Subsection 4.2 describes the published conformance specification. Subsections 4.3, 4.4, and 4.5 present explicit mappings to the IMDA Framework dimensions, the Mao et al. cross-layer attack vector taxonomy, and the Ferrag et al. threat domain taxonomy, respectively.


\subsection{Live Deployment Overview}

The MolTrust reference implementation is live at \texttt{api.moltrust.ch} and has been in production operation since March 2026. Key metrics as of April 2026:


\begin{itemize}
\item \textbf{Eight credential verticals operational:} Core Identity, Commerce, Travel, Skill Verification, Prediction Markets, Brand Protection, Music (via the Moltify subsystem), and Sports Integrity. Each vertical has defined credential schemas, active issuance workflows, and verification endpoints.
\item \textbf{48 Model Context Protocol (MCP) tools} published via the \texttt{@moltrust/openclaw} npm package, compatible with any MCP-enabled AI assistant or development environment.
\item \textbf{44,355 agents indexed} via the ERC-8004 ecosystem scanner \cite{ref37_erc8004}, providing cross-ecosystem agent visibility independent of MolTrust registration.
\item \textbf{Base Layer 2 on-chain anchoring} active for agent registrations, credential issuances, and confirmed violations. Protocol Whitepaper v0.8 and Technical Specification v0.8 are themselves anchored on Base L2 (TechSpec v0.8 at Block 44745864).
\item \textbf{Multi-chain wallet binding} for Ethereum, Base L2, and Solana with chain-native signature schemes (secp256k1 for EVM, Ed25519 for Solana).
\item \textbf{External DID bridging and cross-ecosystem trust score import} with documented weight decay (imported scores at 0.3 relative weight, 45-day half-life versus 90-day for native scores).
\item \textbf{Three enforcement layers operational:} cryptographic (Ed25519 + JCS), API (Trust Score, IPR, revocation), and kernel (Falco eBPF) — all live as of April 2026.
\item \textbf{Published npm packages with passing test suites:} \texttt{@moltrust/sdk} v1.1.0 (14/14 tests passing), \texttt{@moltrust/aae} v1.0.0, \texttt{@moltrust/openclaw} v1.0.1, \texttt{@moltrust/x402} v1.0.1, \texttt{@moltrust/mpp} v1.0.3.
\item \textbf{Payment protocol integration:} active paywall endpoints using x402 and MPP with live USDC payment events recorded on Base L2.
\end{itemize}

The deployment is operated by a small team (one full-time founder, two part-time contributors) using an agent-first operating model: two autonomous GitHub agents (\texttt{MoltyCelBot} and \texttt{VCOne-AI}) handle outreach, onboarding, and documentation contributions under defined cooldown policies. This operational model is itself a demonstration of the protocol's claims — MolTrust's own production operations depend on the trust infrastructure it defines.


\subsection{Published Conformance Specification}

Protocol-level conformance is documented in a published \texttt{CONFORMANCE.md} specification (version 1.0 as of April 2026), maintained at \texttt{github.\allowbreak com/\allowbreak MoltyCel/\allowbreak moltrust-protocol/\allowbreak docs/\allowbreak CONFORMANCE.\allowbreak md} and mirrored at \texttt{moltrust.ch} with a published SHA-256 checksum (\texttt{041c9a5f...}). The conformance specification defines three conformance tiers:


\begin{itemize}
\item \textbf{Layer A (Protocol Standard) conformance} — thirteen requirements covering DID issuance, VC issuance and verification, interaction proof structure, Ed25519Signature2020 over RFC 8785 canonical JSON, UUID v4 identifiers, expiry rejection, delegation chain depth limits, vertical naming format, endorsement format, violation record format, AAE validation rules, and credential TTL enforcement with revocation checking.
\item \textbf{Layer B (Registry) conformance} — eight requirements covering registry API endpoints, trust score response signing, revocation propagation, operator DID publication at \texttt{.well-known/did.json}, violation record association with both agent and principal DIDs, on-chain anchoring format, and Bitstring Status List support.
\item \textbf{Layer C (Reputation Model) conformance} — non-mandatory; implementations MAY use any scoring model consuming Layer A evidence, with optional full implementation of the reference model specified in Technical Specification v0.8 Section 4.
\end{itemize}

Three operational safeguards ensure conformance specification integrity. First, the specification is \textbf{auto-generated} from source code by a deterministic script (\texttt{gen\_conformance.py}), eliminating documentation-implementation drift at generation time. Second, a \textbf{pre-push Git hook} and an hourly \textbf{watchdog process} compare the live checksum against the anchored version; divergence triggers a Telegram alert to the operator. Third, the \textbf{SKILL AUDIT v1.0} runs nine conformance checks at the protocol endpoint level, each with explicit mapping to Common Weakness Enumeration (CWE) identifiers where applicable (Table 1). The SKILL AUDIT endpoints \texttt{/guard/audit/checks} and \texttt{/version} are publicly accessible for third-party conformance verification.




\begin{table}[!htbp]
\centering
\footnotesize
\renewcommand{\arraystretch}{1.15}
\caption{SKILL AUDIT v1.0 conformance checks with CWE mappings}
\begin{tabularx}{\textwidth}{@{}c l X X l@{}}
\toprule
\# & Check & Description & CWE Mapping & Severity \\
\midrule
1 & \texttt{secrets\_scan} & Repository scanned for hard-coded credentials, API keys, private keys & CWE-798 (Use of Hard-coded Credentials) & HARD fail \\
2 & \texttt{tls\_enforcement} & All endpoints served over TLS 1.2+; HSTS enabled & CWE-319 (Cleartext Transmission of Sensitive Information) & HARD fail \\
3 & \texttt{input\_validation} & All JSON inputs validated against RFC 8785 canonical form; oversize inputs rejected & CWE-20 (Improper Input Validation) & Fail \\
4 & \texttt{signature\_verification} & Ed25519 signatures verified against DID Document public keys; no skip paths & CWE-347 (Improper Verification of Cryptographic Signature) & HARD fail \\
5 & \texttt{expiry\_enforcement} & Expired credentials rejected; no grace period; UTC comparison & CWE-613 (Insufficient Session Expiration) & Fail \\
6 & \texttt{rate\_limit} & Rate limits enforced on public endpoints; per-IP and per-API-key & CWE-770 (Allocation of Resources Without Limits or Throttling) & Warning \\
7 & \texttt{cors\_policy} & CORS policies restrict origins; no wildcard for credential endpoints & CWE-346 (Origin Validation Error) & Fail \\
8 & \texttt{dependency\_audit} & npm and pip dependencies scanned for known CVEs; critical vulnerabilities block release & CWE-1104 (Use of Unmaintained Third-Party Components) & Fail \\
9 & \texttt{logging\_integrity} & Operation logs write-protected; critical events signed with operator key & CWE-778 (Insufficient Logging) & Warning \\
\bottomrule
\end{tabularx}
\end{table}

The CWE mappings enable MolTrust conformance checks to cross-reference with NIST National Vulnerability Database entries and with the Ferrag et al. \cite{ref18_ferrag} Protocol Vulnerabilities taxonomy, providing institutional interoperability at the vulnerability-tracking level.


This conformance infrastructure is substantively different from a protocol specification published only in a document. Independent verifiers can reproduce any conformance check, any test vector, and any live endpoint response against a cryptographically anchored specification — without requiring any MolTrust-specific tooling beyond a standard SHA-256 implementation and HTTP client. This is the basis of Novelty Claim B: that protocol-level interoperability is not asserted but demonstrated, through reproducible and independently verifiable evidence.


\subsection{Mapping to Singapore IMDA Framework}

The Singapore IMDA \textit{Model AI Governance Framework for Agentic AI} \cite{ref8_imda} defines four dimensions for governance of autonomous AI systems. Table 2 presents the mapping of MolTrust capabilities to each dimension.




\begin{table}[!htbp]
\centering
\footnotesize
\renewcommand{\arraystretch}{1.15}
\caption{MolTrust mapping to Singapore IMDA MGF for Agentic AI dimensions}
\begin{tabular}{@{}p{2.5cm} p{5.0cm} p{5.0cm} p{2.1cm}@{}}
\toprule
IMDA Dimension & Requirement & MolTrust Mechanism & Status \\
\midrule
1. Risk Bounding & Define limits on agent powers; least-privilege access & AAE MANDATE (purpose enumeration, allowed/denied action patterns, resource patterns) and CONSTRAINTS (financial thresholds, jurisdictions, time windows, counterparty score gates) & Fully Addressed \\
2. Human Accountability & Traceable chains from autonomous action to responsible humans; human-approval checkpoints & Five-Party Trust Chain (Developer $\rightarrow$ Owner $\rightarrow$ Agent $\rightarrow$ Instructor $\rightarrow$ Counterparty); AAE \texttt{requireHumanApprovalAbove} and \texttt{approvalThreshold} fields & Partially Addressed (HITL interface on roadmap Q2 2026) \\
3. Technical Controls & Agent lifecycle identity, monitoring, revocation & Full W3C DID + VC lifecycle (issuance, verification, TTL enforcement, CAEP-compatible revocation); IPRs as continuous audit trail; Falco kernel-layer enforcement & Strongly Addressed \\
4. End-User Responsibility & User understanding of agent capabilities and limitations & Machine-readable AAE metadata and VCs; human-readable Agent Cards on roadmap Q2 2026 & Roadmap Q2 2026 \\
\bottomrule
\end{tabular}
\end{table}

Two dimensions are fully or strongly addressed (1, 3). Two dimensions are partially addressed with specific roadmap commitments (2, 4). This mapping is consistent with the MolTrust design principle that trust infrastructure provides the \textit{verifiable substrate} for accountability — legal accountability and user-experience design are adjacent concerns that must be built on top of the infrastructure, not within it.


\subsection{Coverage of Mao et al. Cross-Layer Attack Vectors}

The Mao et al. Systematization of Knowledge on AI agent security \cite{ref20_mao} identifies twelve cross-layer attack vectors spanning five dimensions of the agent security landscape. Table 3 presents MolTrust's coverage of each vector, with the addressing mechanism and current status.




\begin{table}[!htbp]
\centering
\footnotesize
\renewcommand{\arraystretch}{1.15}
\caption{MolTrust coverage of Mao et al. cross-layer attack vectors}
\begin{tabular}{@{}p{1.0cm} p{2.4cm} p{4.0cm} p{4.3cm} p{2.5cm}@{}}
\toprule
Vector & Name & Description & MolTrust Coverage & Status \\
\midrule
P2T & Principal-to-Task & Malicious or compromised principal issues unauthorized task & AAE VALIDITY block with issuer DID verification; Principal DID continuity; Trust Tier 0 optional KYC anchor & Covered \\
T2R & Task-to-Resource & Task scope exceeds authorized resources & AAE MANDATE.resources with ABAC semantics; default-deny resource matching & Covered \\
A2M & Agent-to-Model & Agent attempts unauthorized model interaction & AAE MANDATE.allowedActions with URI-pattern matching; \texttt{toolAllowlist} in CONSTRAINTS.obligations & Covered \\
R2I & Resource-to-Interaction & Transactions chain to produce aggregate unauthorized state (e.g., wash trading, layering) & AML pattern detection & Roadmap Phase 2/3 (3–4 weeks) \\
T2T & Task-to-Task & Task chain violates separation of duties & AAE delegation semantics (attenuation-only, maximum depth 8, deny-precedence) & Covered \\
P2K & Principal-to-Key & Principal key compromise & Key rotation with revoked-key retention in DID Document; CAEP-compatible revocation; Violation Records with principal DID association & Covered \\
M2A & Model-to-Agent & Model integrity compromise (backdoor, poisoning) & Partnership space (Sigstore, external model-integrity attestation) & Explicit Non-Coverage \\
C2E & Context-to-Execution & Context poisoning leads to unauthorized execution & Layer 3 Falco kernel enforcement (execution-time detection independent of agent runtime) & Covered \\
O2P & Oracle-to-Protocol & Oracle manipulation produces false ground-truth & Chainlink / oracle-layer space & Explicit Non-Coverage \\
I2M & Identity-to-Market & Sybil identities manipulate market-level signals & Dual-signature IPRs, cross-vertical diversity requirement, Jaccard cluster detection (threshold 0.8), x402 economic cost per interaction, permanent Violation Records & Covered \\
N2C & Network-to-Capital & Network-level transactions produce cumulative capital exposure (structuring detection) & Cumulative exposure monitoring & Roadmap Phase 2 (1–2 weeks post-current milestone) \\
S2I & Settlement-to-Identity & Settlement artifacts enable identity de-anonymization & Outcome hash construction with entropy requirements; separation of on-chain anchors (hashes only) from off-chain raw data & Covered \\
\bottomrule
\end{tabular}
\end{table}

\textbf{Coverage summary:} Eight of twelve vectors are addressed by live capabilities (P2T, T2R, A2M, T2T, P2K, C2E, I2M, S2I). Two vectors are explicit roadmap items with published milestones (R2I, N2C). Two vectors are explicit non-coverage: M2A requires a model-integrity attestation partnership (Sigstore is the current reference target); O2P is outside the scope of a trust infrastructure layer and belongs to the oracle-layer ecosystem (Chainlink and equivalent). This distinction — \textit{covered}, \textit{roadmap}, \textit{explicit non-coverage} — is significant: it resists the temptation to overclaim coverage by counting partial or theoretical mitigations, and it provides a machine-readable commitment schedule for gaps that will close within a defined timeframe.


\subsection{Mapping to Ferrag et al. Threat Domains}

Ferrag et al. \cite{ref18_ferrag} organize LLM-agent threats into four domains. MolTrust's relationship to each domain is the following.


\textbf{Input Manipulation} (prompt injections, long-context hijacks, multimodal adversarial inputs) is \textit{outside MolTrust's scope}. Input manipulation targets the agent's reasoning layer; trust infrastructure operates beneath the reasoning layer. MolTrust's contribution to input-manipulation defense is indirect: an agent that executes a prompt-injection-induced unauthorized action still leaves a signed AAE trail, and the resulting unauthorized action — if detected — can be recorded as a Violation Record that propagates to the agent's trust score. MolTrust does not prevent input manipulation; it ensures that the \textit{consequences} of successful manipulation are observable and accountable.


\textbf{Model Compromise} (backdoors, poisoning, prompt-parameter attacks) is also \textit{outside MolTrust's scope}. This domain is addressed by model-integrity attestation frameworks (Sigstore, TEE-based evaluation, verifier-independent safety benchmarks such as those proposed by Schnabl et al. \cite{ref33_schnabl}). MolTrust's M2A position (Section 4.4) is explicit non-coverage.


\textbf{System and Privacy Attacks} (side-channels, membership inference, retrieval poisoning) are \textit{partially in scope}. MolTrust addresses retrieval-provenance attacks through the Interaction Proof Record mechanism (every credential or record consumed by an agent has a verifiable issuer chain). Side-channel attacks on the MolTrust infrastructure itself are addressed by operational controls (TLS, rate limiting, input validation) outside the scope of this paper. Membership inference on credential issuance is partially mitigated by the explicit public-verifiability design of the protocol: credential issuance is by design observable, making membership inference non-adversarial for registered agents.


\textbf{Protocol Vulnerabilities} (MCP, ACP, ANP, A2A exploits) are \textit{directly in scope} and are the primary domain where MolTrust's contribution is concrete. Each of the protocol layers Ferrag et al. identify — MCP, A2A — is addressed through MolTrust's published SDK and middleware: \texttt{@moltrust/sdk} for HTTP transport with the \texttt{middleware()}/\texttt{register()}/\texttt{verify()} API; \texttt{@moltrust/openclaw} for MCP server binding; active integration threads in the A2A project repository. The Skill Audit infrastructure implements nine protocol-level conformance checks at endpoint boundaries, each mapped to CWE identifiers where applicable. This is the domain where the CVE/NIST NVD cross-mapping methodology Ferrag et al. establish provides the most direct reference.


The coverage pattern is consistent across the three threat taxonomies: MolTrust addresses the \textit{infrastructure} layer of agent security — identity, authorization, behavioral record — and does not claim to address the \textit{model} or \textit{input} layers. This scope discipline is deliberate and is consistent with the framework convergence Section 2.2 identifies: Anthropic, Google, and Microsoft each acknowledge that the model and input layers require separate, complementary mitigations. MolTrust fills one specific gap in the overall agent security stack. It does not fill all of them.



\section{Cross-Protocol Interoperability}

This section presents the empirical basis for Novelty Claim B: that MolTrust's protocol-level interoperability is demonstrated rather than asserted, through published specifications, reproducible test vectors, and active integration with independent verifier implementations.


\subsection{qntm Authority Constraints Mapping}

The qntm Authority Constraints specification defines a ConstraintEvaluation model with five facets: \texttt{scope}, \texttt{spend}, \texttt{time}, \texttt{reputation}, and \texttt{reversibility}. The MolTrust AAE maps directly to each facet:


\begin{itemize}
\item \texttt{mandate.purpose} and \texttt{mandate.allowedActions} $\rightarrow$ qntm \texttt{scope} (URI patterns map directly to qntm scope constraints)
\item \texttt{constraints.\allowbreak limits.\allowbreak autonomousThreshold} $\rightarrow$ qntm \texttt{spend} (USDC threshold maps directly to qntm spend limit)
\item \texttt{constraints.\allowbreak duration.\allowbreak ttl} $\rightarrow$ qntm \texttt{time} (seconds TTL maps directly to qntm time constraint)
\item \texttt{constraints.\allowbreak scope.\allowbreak counterpartyMinScore} $\rightarrow$ qntm \texttt{reputation} (0–100 score maps directly to qntm reputation threshold)
\item \texttt{mandate.\allowbreak delegation.\allowbreak allowed} $\rightarrow$ qntm \texttt{reversibility} (inverse mapping: \texttt{delegation=false} implies non-reversible commitment)
\end{itemize}

The mapping was verified through five conformance test vectors (TV-001 through TV-005) published in the \texttt{corpollc/qntm} repository pull request \#10 (https://github.com/corpollc/qntm/pull/10). Each test vector exercises AAE delegation narrowing semantics: TV-001 covers top-level agent delegation; TV-002 and TV-003 cover sub-agent delegation at depths 2 and 3; TV-004 covers deny-precedence semantics (an action matching both \texttt{allowedActions} and \texttt{deniedActions} must be denied); TV-005 covers attenuation enforcement (a sub-agent scope exceeding parent scope must be rejected). All five test vectors verify against both the MolTrust reference implementation and the qntm ConstraintEvaluation verifier, using shared primitives: JCS RFC 8785 canonicalization and Ed25519 signatures.


The structural consequence is that a single signing operation can produce artifacts verifiable by both MolTrust and qntm conformant implementations, without format translation. This is not interoperability by specification — it is interoperability by shared primitive. The five test vectors are reproducible by any third party against the live \texttt{api.moltrust.ch} endpoint using standard tooling (\texttt{openssl}, \texttt{jose}, or any Ed25519 library with RFC 8785 support).


A scope note on the qntm mapping: qntm ConstraintEvaluation is designed with Datalog-based policy semantics derived from the Biscuit token format \cite{ref34_biscuit}. MolTrust's AAE uses URI-pattern matching for action and resource specification, which is simpler to implement and audit but less expressive for complex conditional authorization (nested logical constraints, multi-party temporal conditions, recursive rules). Datalog-level policy expressiveness is a documented MolTrust roadmap item for Q3–Q4 2026; the current URI-pattern approach is acknowledged as an expressiveness limitation in the companion AIP implementation paper \cite{ref27_kroehlssrn}.


\subsection{APS Provider Attestation Mapping}

The Agent Provider Service (APS) protocol uses a provider-attestation model for agent credential exchange. The APS \texttt{importProviderAttestation()} entry point accepts MolTrust W3C Verifiable Credentials as attestation input. The MolTrust Trust Score (0–100) maps to APS Grade levels via a defined quantization: Grade 0 = score less than 25; Grade 1 = 25 to less than 50; Grade 2 = 50 to less than 75; Grade 3 = 75 or above.


The APS Grade mapping is unidirectional: MolTrust credentials import into APS through the provider-attestation interface, producing an APS-consumable Grade attestation. The reverse direction — APS credentials importing into MolTrust — is conceptually supported through the external DID bridging mechanism (Section 3.1 of \cite{ref28_techspec}) and external trust score import, but is not currently an active integration path as APS is itself in early deployment.


An analysis of the APS \texttt{aps.txt} provider-discovery mechanism identified five attack vectors (DNS spoofing, content manipulation, replay, denial-of-service, and conflicting-policy confusion); MolTrust contributed a signed-aps.txt proposal as a Working Group contribution, with mitigation mapping for each vector. This contribution is documented in [28, Section 10.2] and in the OWASP LLM Top 10 Runtime Integrity Layers discussion thread [19, Issue \#802].


\subsection{Payment Protocol Composability}

MolTrust is designed as a payment-protocol-agnostic trust layer: the agent registers a W3C DID with MolTrust and receives a Trust Score queryable by any endpoint operator, regardless of which payment rail the agent uses. This is the structural advantage of building on W3C DID and VC standards rather than protocol-specific identity systems — portability is guaranteed by the standard, not by bilateral agreements. Current payment protocol integrations:


\begin{itemize}
\item \textbf{x402} (Coinbase, Cloudflare): \texttt{@moltrust/x402} v1.0.1 published; HTTP 402 challenge-response middleware with \texttt{requireScore(\{ minScore, failBehavior \})} API; live endpoints behind active x402 paywall.
\item \textbf{MPP — Machine Payments Protocol} (Stripe, Tempo, Visa): \texttt{@moltrust/mpp} v1.0.3 published; session-based payments, Stripe Single-use Payment Token support, Bitcoin Lightning support; identical \texttt{requireScore()} API to \texttt{@moltrust/x402} for multi-protocol endpoint operators.
\item \textbf{AP2 — Agent Payments Protocol} (Google, 60+ partners): planned for MolTrust Protocol v0.9 (Q3 2026).
\item \textbf{ACP — Agent Commerce Protocol} (OpenAI, Stripe): planned for MolTrust Protocol v0.9 (Q4 2026).
\end{itemize}

The shared \texttt{requireScore()} API across the live integrations is deliberate: it means an endpoint operator deploying MolTrust middleware does not need separate integration code per payment protocol. The trust gate is unified; the payment rail is fungible. This is the operational realization of the composability claim Section 2.2 identifies as Anthropic's framework objective.


\subsection{Additional Vocabulary Integration Work}

Beyond the qntm and APS integrations that anchor Novelty Claim B, MolTrust maintains additional vocabulary and registration integration work with adjacent agent-governance projects, including an active proposal thread in the OpenClaw project (RFC \#49971 \cite{ref35_openclaw} on native agent identity integration), contributions to multiple agent-governance vocabulary efforts, cross-verification with independent crosswalk templates, and ERC-8004 agent registration. These integrations are at earlier maturity stages than the qntm and APS mappings and are documented here for completeness rather than as primary evidence for the interoperability claim. The qntm test vectors (TV-001 through TV-005) and the APS Grade mapping remain the reference implementations for the cross-protocol interoperability argument.



\section{Sybil Resistance and Behavioral Consistency}

This section details the Sybil resistance mechanisms that constitute Claim A(c): layered defenses combining cryptographic, structural, and economic barriers against identity fabrication and endorsement manipulation. The mechanisms described here operate within the trust score architecture defined in Section 3.1 and are designed to compound — each mechanism raises the cost of a Sybil campaign independently, and their combination raises the cost multiplicatively.


\subsection{Three-Mechanism Sybil Resistance}

MolTrust's Sybil resistance is organized around three mechanisms, each addressing a distinct attack surface.


\textbf{Mechanism 1: Dual-Signature Interaction Proofs.} Interaction Proof Records (IPRs) require cryptographic signatures from two distinct Ed25519 signing keys — one from the initiator and one from the counterparty. The responder's signature covers the initiator's signature, creating a sequential commitment chain that is non-repudiable by either party. Forging an IPR requires simultaneous compromise of two independent keypairs, which transforms the Sybil attack from an identity-fabrication problem (cheap) to a key-compromise problem (expensive). A single adversary controlling multiple agent DIDs can produce self-endorsements, but cannot produce bilateral IPRs without a cooperating counterparty — and each cooperating counterparty is itself subject to endorsement-pattern analysis.


The dual-signature requirement serves a second function beyond Sybil resistance: it produces an auditable record of actual interactions between identified parties. Even if a Sybil campaign successfully inflates an agent's endorsement count, the absence of corresponding bilateral IPRs from independent counterparties is itself a detectable signal. The Trust Score formula (Section 3.1) weights the interaction bonus separately from endorsement-derived scores, creating a structural divergence between endorsement-inflated agents and genuinely active agents.


\textbf{Mechanism 2: Cross-Vertical Diversity Gate.} As described in Section 3.1, cross-vertical diversity operates as a two-stage mechanism. The incentive layer rewards breadth of endorsement across the seven canonical endorsement verticals (\texttt{core}, \texttt{skill}, \texttt{shopping}, \texttt{travel}, \texttt{prediction}, \texttt{salesguard}, \texttt{sports}) via a weighted bonus capped at 30 points. The gate layer applies a flat Sybil penalty of 10.0 when a non-seed agent has endorsements from fewer than three distinct verticals, producing effective score nullification after the $\times$20 multiplier (10.0 $\times$ 20 = 200, clamped to [0, 100]).


The coordination-cost argument is direct: two colluding agents can fabricate mutual endorsements within a single vertical at minimal cost. Requiring endorsements from three independent verticals forces the Sybil operator to maintain distinct credibility signals across unrelated domains — skill verification, commerce, travel — each with its own credential schemas and verification workflows. The cost of maintaining credible presence across three verticals scales non-linearly with the number of Sybil identities, as each vertical requires independent credential issuance and verification interactions.


\textbf{Mechanism 3: Principal-DID-Linked Violation Persistence.} Violation records — including MoltGuard security findings, Sybil detection flags, and repetitive endorsement pattern alerts — are associated with the principal DID, not the agent DID. When an agent is revoked and re-registered under a new DID, the principal's violation history persists. An agent operator cannot rotate identity to escape a negative trust history; the Five-Party Trust Chain (Section 3.2) ensures that the Developer and Owner DIDs carry the accumulated behavioral record across agent re-registrations.


This mechanism addresses the identity-rotation attack that is the simplest and most common Sybil strategy in permissionless systems: create a new identity when the old one accumulates negative signals. In MolTrust, creating a new agent DID under the same principal DID inherits the principal's violation history. Creating a new principal DID requires a new Trust Tier 0 KYC credential (where applicable) or begins with zero trust history and must re-accumulate endorsements through the three-vertical gate — which returns the attacker to the coordination-cost barrier of Mechanism 2.


\subsection{Jaccard Cluster Detection Mechanics}

The Jaccard similarity coefficient measures endorser-set overlap between agents. For two agents \textit{i} and \textit{j} with endorser sets \textit{E\_i} and \textit{E\_j}, the Jaccard index is |\textit{E\_i} $\cap$ \textit{E\_j}| / |\textit{E\_i} $\cup$ \textit{E\_j}|. When this index exceeds the production threshold of 0.8, a Sybil penalty is applied proportional to the similarity and the endorser count:


\begin{lstlisting}[basicstyle=\ttfamily\small]
penalty = jaccard x num_endorsers x 0.5
\end{lstlisting}

The penalty feeds into the final score formula as \texttt{sybil\_penalty $\times$ 20}. A worked example: Agent X has four endorsers {A, B, C, D}. Agent Y has endorsers {A, B, C, E}. The Jaccard index is |{A,B,C}| / |{A,B,C,D,E}| = 3/5 = 0.6, which is below the 0.8 threshold — no penalty. Now consider a Sybil cluster: Agent X has endorsers {A, B, C, D} and Agent Z has endorsers {A, B, C, D} (identical sets). Jaccard = 4/4 = 1.0, exceeding 0.8. Penalty = 1.0 $\times$ 4 $\times$ 0.5 = 2.0. After the $\times$20 multiplier: 40-point score reduction. For a moderately-endorsed agent with a raw score of 65, this reduces the final score to 25 — effectively dropping the agent below the threshold required for commercial interactions.


The 0.8 threshold is calibrated to allow legitimate endorser overlap (agents in the same ecosystem may share some endorsers) while penalizing near-identical endorser sets that indicate coordinated fabrication. The threshold is a configuration parameter; formal sensitivity analysis across the [0.6, 0.95] range is planned once the agent population exceeds 200 registered agents, providing sufficient statistical power for meaningful comparison.


\subsection{Economic Cost Layer}

The economic cost layer functions as an anti-spam compounding factor rather than a standalone Sybil defense. While x402 per-verify pricing adds negligible friction for small-scale Sybil campaigns, it compounds with the cryptographic and structural mechanisms of Sections 6.1–6.2: a Sybil operator must simultaneously forge dual signatures, distribute endorsements across three independent verticals, and absorb per-interaction economic cost at scale. The compounding effect, not any single layer, is the defense.


At current pricing (agent score query at \$0.05, sybil scan at \$0.10), a campaign requiring 1,000 score verifications to establish Sybil credibility incurs \textasciitilde{}\$50–100 in direct payment costs, plus Base L2 gas fees for each on-chain payment event. The economic barrier is most effective against low-effort, high-volume Sybil strategies — the "spray and pray" attacks that dominate permissionless systems — rather than against targeted, well-resourced adversaries.


\subsection{Trust Score Weights: Empirical Selection}

The weight assignment ($\alpha$ = 0.6, $\beta$ = 0.3, $\gamma$ = 0.1) was selected empirically during the bootstrap phase of the registry. The 2:1 ratio between direct and propagated scores reflects a design preference for direct evidence over transitively-inherited trust: an endorsement from an agent with whom the subject has interacted directly is weighted twice as heavily as second-order reputation propagated through the endorsement graph. The comparatively small $\gamma$ = 0.1 weight on the cross-vertical bonus is intentional; the primary cross-vertical defense is the hard gate (Section 3.1, Mechanism 2 above), not the bonus. The bonus rewards breadth; the gate enforces a minimum.


These weights are preliminary. A formal sensitivity analysis across alternative weightings (e.g., 0.5/0.4/0.1, 0.7/0.2/0.1, equal-weight 0.33/0.33/0.33) is planned once the agent population exceeds 200 registered agents with sufficient endorsement density for statistical comparison of score distributions. The current weights produce a distribution that — at the present scale of 54 agents and 57 endorsements — separates seed agents, actively-endorsed agents, and minimally-endorsed agents into distinct score bands. Whether this separation is maintained under adversarial load at 10$\times$ or 100$\times$ scale is an open empirical question addressed in Section 7.2.


\subsection{Scale Gap}

No adversarial Sybil campaigns have been observed or deliberately simulated against the live registry during the two-month production window (54 agents, 57 endorsements as of April 2026). The three-mechanism design is theoretically motivated by Sybil coordination-cost arguments and operates at the registry's current scale without observed bypasses, but empirical validation under adversarial load — including red-team Sybil campaigns with 100 or more colluding agents — has not been conducted.


The relevant external benchmark is Agent.market's reported scale of 69,000 autonomous bots \cite{ref1_agentmarket}. MolTrust's three-mechanism resistance is designed with that scale in mind: dual-signature IPRs scale linearly with interaction count, Jaccard detection scales quadratically with endorser-set size (each new agent is compared against existing agents), and the cross-vertical gate is O(1) per agent. Computational scalability is not the concern; the concern is whether the heuristic thresholds (Jaccard 0.8, three-vertical minimum, penalty multiplier $\times$20) remain effective when a sophisticated adversary with resources to maintain hundreds of credible identities across multiple verticals targets the registry specifically. This is treated as a primary limitation in Section 7.2.


\subsection{Trust Score Formula Reference}

The trust score computation is fully specified in Section 3.1 and normatively defined in Section 4 of the Technical Specification \cite{ref28_techspec}. The structural components are: four weighted inputs (direct-endorsement score at $\alpha$ = 0.6, propagated endorsement score at $\beta$ = 0.3, cross-vertical diversity bonus at $\gamma$ = 0.1, additive interaction bonus capped at 10 points); two modifiers (Sybil penalty at $\times$20 multiplier, inactivity penalty as additive adjustment); and the seed-floor-guard ensuring seed agents do not fall below their registered base\_score. The final score is clamped to [0, 100] and withheld for non-seed agents with fewer than three distinct endorsers.



\section{Discussion and Limitations}

\subsection{What the Three Claim A Capabilities Provide}

The three capabilities constituting Claim A — kernel-layer AAE enforcement, cross-protocol interoperability, and layered Sybil resistance — address distinct layers of the agent trust problem and are independently verifiable.


Kernel-layer AAE enforcement via Falco eBPF (Section 3.4) provides behavioral integrity monitoring below the agent process boundary. This capability addresses the specific gap that Rodriguez Garzon et al. \cite{ref36_rodriguezgarzon} identify when an agent's LLM is "in sole charge of security procedures" — the enforcement mechanism operates independently of the agent's own runtime and cannot be suppressed from userspace. It is the structural response to the "policy self-modification between authorization and execution" problem.


Cross-protocol interoperability (Section 5) is demonstrated through five published test vectors (TV-001 through TV-005) verified against independent verifier implementations: qntm Authority Constraints ConstraintEvaluation and APS Provider Attestation. The structural consequence is that a single MolTrust signing operation produces artifacts verifiable by conformant implementations of multiple independent protocols, without format translation. This is interoperability by shared cryptographic primitive (Ed25519 + RFC 8785 JCS), not by bilateral specification agreement.


Layered Sybil resistance (Section 6) combines three mechanisms — dual-signature IPRs, cross-vertical diversity gating, and principal-DID-linked violation persistence — that compound coordination costs for identity-fabrication attacks. Each mechanism is independently deployable and independently verifiable; their combination raises the minimum viable Sybil campaign cost beyond casual or automated attack strategies.


These three capabilities position MolTrust relative to adjacent academic work. SAGA \cite{ref16_syros} provides centralized policy enforcement with formal security guarantees but requires a trusted Provider — it cannot operate across organizational boundaries where no such Provider exists. Acharya's TIVA \cite{ref17_acharya} describes the architectural pattern (DIDs + VCs + on-chain anchoring) as a theoretical framework with a conceptual architecture diagram but without a deployed reference implementation or published conformance specifications. Rodriguez Garzon et al. \cite{ref36_rodriguezgarzon} demonstrate cross-domain trust establishment between DID-equipped agents and explicitly identify the limitation of LLM-only governance, but do not implement enforcement below the model layer. MolTrust contributes the deployed realization: the same architectural primitives, with kernel-layer enforcement, published conformance evidence, and reproducible test vectors against independent verifier implementations.


\subsection{Limitations}

\subsubsection{7.2.1 Scale Discipline}

The MolTrust reference registry operates with 54 registered agents, 57 endorsements, and a two-month production history (March–April 2026). These numbers are honest: they represent a functioning production system, not a simulation or testbed, but they do not constitute validation at the scale of the agent economy the introduction describes. The three-mechanism Sybil resistance has not been tested under adversarial load, and the trust score weight calibration has not been validated against a statistically significant agent population. Claims in this paper are bounded by the evidence available at this scale.


\subsubsection{7.2.2 Heuristic Nature of Sybil Detection}

The Jaccard similarity threshold (0.8), the three-vertical minimum gate, and the $\times$20 penalty multiplier are empirically calibrated parameters, not formally proven optimal values. The Jaccard threshold tolerates legitimate endorser overlap up to 80\% while penalizing near-identical endorser sets; the three-vertical gate reflects a coordination-cost argument rather than a formal security proof. Alternative threshold values have not been systematically compared. Formal sensitivity analysis is planned once the registry reaches sufficient scale (N $\geq$ 200 agents) for meaningful distributional comparison. Until then, these parameters should be understood as operational heuristics informed by the bootstrap deployment, not as security guarantees.


\subsubsection{7.2.3 URI-Pattern vs. Datalog Expressiveness}

The AAE constraint model uses URI-pattern matching for action and resource specification. This is simpler to implement, audit, and explain than full Datalog or Relationship-Based Access Control (ReBAC) expressiveness, but it limits the protocol's ability to express complex conditional authorization: nested logical constraints, multi-party temporal conditions, recursive delegation rules. For the current deployment — where authorization boundaries are defined by endpoint operators with straightforward allow/deny patterns — URI-pattern matching is sufficient. For agent-to-agent delegation graphs with many levels of conditional sub-delegation, Datalog-level expressiveness would be required. This is a documented roadmap item for Protocol v0.9 (Q3–Q4 2026); the current approach is acknowledged as an intentional trade-off favoring deployability over expressiveness.


\subsubsection{7.2.4 eBPF/Falco Bypass Scenarios}

Kernel-layer enforcement via Falco eBPF is not a complete isolation guarantee. Known bypass scenarios include: (i) agents running inside privileged containers that share the host kernel namespace and can disable or reconfigure eBPF probes; (ii) container-escape vulnerabilities (CVE-2019-5736 class) that compromise the host kernel itself; (iii) rule-evasion via syscall sequences not covered by the current Falco rule set; (iv) eBPF probe disablement by operator-level access to the container runtime. MolTrust's kernel-layer enforcement assumes non-privileged agent containers running on a non-compromised host with an operator-maintained Falco rule set. Operators running agents in privileged mode must be aware that kernel-layer enforcement is advisory, not enforced, in that deployment configuration. Layer 3 is a defense-in-depth complement to Layers 1 and 2, not a standalone security boundary.


\subsubsection{7.2.5 No Third-Party Security Audit}

An internal security review conducted in March 2026 identified and resolved findings across CORS configuration, hardcoded credential exposure, SVG-based XSS vectors, CSRF token handling, and race conditions in credential issuance. All identified findings were remediated prior to production deployment. No third-party penetration test or formal security audit has been commissioned as of publication. The SKILL AUDIT v1.0 (Table 1, Section 4.2) provides automated conformance checking with CWE-mapped vulnerability coverage, but this is not a substitute for independent adversarial security assessment. A third-party audit is on the post-funding roadmap and is identified as a prerequisite for enterprise adoption.


\subsubsection{7.2.6 Two-Month Production Horizon}

The two-month operational window (March–April 2026) limits claims about long-term trust dynamics. Trust score decay behavior over periods exceeding the 90-day half-life has not been observed empirically. Reputation inheritance patterns across agent re-registrations have been implemented but not exercised at meaningful scale. Adversarial adaptation cycles — where attackers observe the detection heuristics and adjust their strategies — have not been observed because no adversarial campaigns have targeted the registry. Long-term trust dynamics are inherently difficult to validate in a system that has been operational for two months; this paper reports what has been built and deployed, not what has been validated over time.


\subsection{Positioning Against Adjacent Work}

The MolTrust Protocol is positioned relative to four categories of adjacent work.


Against centralized architectures (SAGA \cite{ref16_syros}), MolTrust trades the formal security guarantees of a trusted Provider for cross-organizational portability. SAGA's One-Time Key issuance and policy enforcement are tighter within a single administrative domain; MolTrust's DID-based verification is weaker in policy enforcement granularity but does not require participants to trust a common Provider. The two architectures address different deployment scenarios and are complementary rather than competing.


Against theoretical frameworks (Acharya TIVA \cite{ref17_acharya}), MolTrust provides the deployed reference implementation that validates the architectural pattern. TIVA describes DIDs + VCs + on-chain anchoring as a conceptual architecture; MolTrust demonstrates that this architecture operates in production with published conformance specifications, reproducible test vectors, and kernel-layer enforcement. The contribution is empirical validation, not architectural novelty.


Against conceptual DID/VC frameworks for agents (Rodriguez Garzon et al. \cite{ref36_rodriguezgarzon}), MolTrust extends beyond the trust-establishment dialogue to include authorization enforcement and behavioral monitoring. Rodriguez Garzon et al. explicitly identify the limitation of LLM-only governance of security procedures; MolTrust's Layer 3 kernel enforcement is a direct response to that finding, operating below the model layer entirely.


Against on-chain agent registries (ERC-8004 \cite{ref37_erc8004}), MolTrust adds the credential, authorization, and behavioral layers that on-chain registration alone does not provide. ERC-8004 establishes that an agent identity exists on-chain; MolTrust establishes what that agent is permitted to do, whether it has acted within those permissions, and what its behavioral record looks like across time. The MolTrust ERC-8004 scanner (indexing 44,355 agents) demonstrates that the two systems are complementary: on-chain registration provides discoverability; off-chain trust infrastructure provides accountability.


\subsection{Future Work}

Six concrete items constitute the near-term and mid-term roadmap:


\textbf{N2C Cumulative Exposure Monitoring (Phase 2, near-term).} Detection of structuring patterns where multiple small transactions aggregate to circumvent per-transaction thresholds. Implementation is scoped and estimated at 1–2 weeks post-current milestone.


\textbf{R2I AML Pattern Detection (Phase 2/3, mid-term).} Identification of wash-trading, layering, and circular-transaction patterns across the interaction graph. Requires endorsement density above current levels for meaningful pattern detection.


\textbf{W3C DID Method Registry Submission.} The \texttt{did:moltrust} method specification is prepared for submission to the W3C DID Method Registry. This submission subjects the method to external specification review and is planned for Protocol v0.9.


\textbf{AP2 and ACP Protocol Integrations.} Integration with Google's Agent Payments Protocol (AP2) and OpenAI's Agent Commerce Protocol (ACP) is planned for Protocol v0.9 when these specifications stabilize sufficiently for conformance testing.


\textbf{Third-Party Security Audit.} Independent penetration testing and formal security audit, identified as a prerequisite for enterprise adoption and planned for the post-funding phase.


\textbf{Formal Sensitivity Analysis of Trust Score Weights.} Systematic comparison of alternative weight configurations ($\alpha$, $\beta$, $\gamma$) across the score distribution, planned when the registry reaches N $\geq$ 200 agents with sufficient endorsement density for statistical validity.



\section{Conclusion}

This paper has presented MolTrust, a production-deployed trust infrastructure for autonomous AI agents built on W3C Verifiable Credentials 2.0 and Decentralized Identifiers v1.0, with on-chain anchoring on Base Layer 2.


The paper documents three distinguishing capabilities constituting Claim A. Kernel-layer Agent Authorization Envelope enforcement via Falco eBPF integration operates below the agent process boundary, providing behavioral integrity monitoring that the agent's own runtime cannot suppress or modify — addressing the "policy self-modification between authorization and execution" gap that industry frameworks identify as the hardest open problem in agent identity. Cross-protocol interoperability is demonstrated through five published test vectors (TV-001 through TV-005) verified against independent verifier implementations (qntm Authority Constraints, APS Provider Attestation), establishing that a single MolTrust signing operation produces artifacts verifiable by multiple independent protocol implementations without format translation (Claim B). Layered Sybil resistance combines dual-signature Interaction Proof Records, cross-vertical endorsement diversity gating with a three-vertical minimum threshold, and principal-DID-linked violation persistence that survives agent re-registration — compounding coordination costs for identity-fabrication attacks across cryptographic, structural, and economic dimensions.


The published conformance specification (CONFORMANCE.md v1.0) with SHA-256 drift detection, nine CWE-mapped SKILL AUDIT checks (Table 1), and explicit coverage mappings to three independent threat taxonomies (Singapore IMDA Framework, Mao et al. cross-layer attack vectors, Ferrag et al. threat domains) provide independently reproducible evidence of the protocol's security posture and regulatory alignment.


MolTrust is a trust-enforcement layer for autonomous AI agents. It is not a replacement for identity and access management systems, not a universal agent framework, and not a model-layer safety mechanism. It addresses one specific gap in the agent security stack — the infrastructure layer comprising identity, authorization, and behavioral accountability — and does so using open, W3C-standardized primitives that are verifiable by any party independently of the issuer's platform.


The limitations are explicit: 54 agents, 57 endorsements, two months of production operation, no adversarial-scale validation, no third-party security audit, heuristic Sybil detection parameters not formally optimized. These limitations bound the claims in this paper to what has been built and deployed, not what has been validated at scale.


The contribution is deployment-first evidence. Acharya \cite{ref17_acharya} described the DID + VC + on-chain architecture as a theoretical framework. Rodriguez Garzon et al. \cite{ref36_rodriguezgarzon} demonstrated cross-domain trust establishment between DID-equipped agents and identified the limitation of LLM-only governance. MolTrust demonstrates that this class of architecture operates in production today, with kernel-layer enforcement that addresses the governance limitation Rodriguez Garzon et al. identify, with cross-protocol conformance evidence verified against independent implementations, and with layered Sybil resistance operating across cryptographic, structural, and economic dimensions. The trust infrastructure that regulators, industry, and academia have converged on is not a future requirement. It is implementable now, with existing standards, on existing infrastructure.


